% GENERATED FILE. Edit canonical lessons, then run npm run generate:books.
% canonical-source-hash: fnv1a64:68b40e115bdb89f6
% canonical-lessons: FR-C526-machine-a-laver, FR-C526-essorer, FR-C526-etendre, FR-C526-pince-a-linge, FR-C526-cintre

\chapter{Laver le linge}
\label{ch:fr-laver-le-linge}
\hlchaptermodality{526}

\begin{chapteropening}
\textbf{By the end of this chapter:} \emph{I can talk about doing the washing.}

Complete the last lesson of chapter 526: I can talk about doing the washing.
\end{chapteropening}

\section[la machine à laver]{la machine à laver — washing machine}

\label{lesson:FR-C526-machine-a-laver}



\begin{quote}
Before the new one: say the French for documentary, then the French for to repeat.
\end{quote}

\subsection*{You'll want to know: la machine à laver}
\textbf{la machine à laver} — \textquotedblleft{}washing machine\textquotedblright{}.

\begin{quote}
La machine à laver tourne depuis une heure.
\end{quote}

\begin{quote}
\emph{The washing machine has been running for an hour.}
\end{quote}

Also \emph{le lave-linge}. \emph{Faire une machine} — to do a load of washing.

\begin{cousinweb}
\emph{la machine à laver} joins \emph{machine}, a learned borrowing from Latin \emph{māchina} \textquotedblleft{}a device\textquotedblright{}, and \emph{laver} \textquotedblleft{}to wash\textquotedblright{}, inherited from Latin \emph{lavāre} (Lewis \& Short).
\end{cousinweb}

\begin{grammarlens}[title={What you've built}]
One more word for telling, explaining and describing at B1.
\end{grammarlens}

\subsection*{Your turn}
\begin{itemize}
\raggedright
  \item \emph{Say it:} \emph{la machine à laver}
  \item \emph{Say it:} \emph{la machine à laver}, once more
  \item \emph{Say it:} a sentence of your own with \emph{la machine à laver}
  \item \emph{Recall:} say the French for documentary, then the French for to repeat, then say \emph{la machine à laver} again
\end{itemize}

\begin{tcolorbox}[breakable,colback=teal!4,colframe=teal!35!black,title={Before you move on}]
What does la machine à laver mean? (\textquotedblleft{}Washing machine\textquotedblright{}.) Say it once more.
\end{tcolorbox}
\section[essorer]{essorer — to spin (dry), to wring}

\label{lesson:FR-C526-essorer}



\begin{quote}
Before the new one: say the French for to repeat, then the French for washing machine.
\end{quote}

\subsection*{You'll want to know: essorer}
\textbf{essorer} — \textquotedblleft{}to spin (dry), to wring\textquotedblright{}.

\begin{quote}
La machine essore le linge à la fin du programme.
\end{quote}

\begin{quote}
\emph{The machine spins the washing at the end of the cycle.}
\end{quote}

\emph{L'essorage} — the spin cycle.

\begin{cousinweb}
\emph{essorer} comes from Old French \emph{essorer} (EtymDB), from Vulgar Latin \emph{exaurāre}, built on Latin \emph{aura} \textquotedblleft{}a breeze\textquotedblright{} (Lewis \& Short): to dry washing was first to put it out in the air.
\end{cousinweb}

\begin{grammarlens}[title={What you've built}]
One more word for telling, explaining and describing at B1.
\end{grammarlens}

\subsection*{Your turn}
\begin{itemize}
\raggedright
  \item \emph{Say it:} \emph{essorer}
  \item \emph{Say it:} \emph{essorer}, once more
  \item \emph{Say it:} a sentence of your own with \emph{essorer}
  \item \emph{Recall:} say the French for to repeat, then the French for washing machine, then say \emph{essorer} again
\end{itemize}

\begin{tcolorbox}[breakable,colback=teal!4,colframe=teal!35!black,title={Before you move on}]
What does essorer mean? (\textquotedblleft{}To spin (dry), to wring\textquotedblright{}.) Say it once more.
\end{tcolorbox}
\section[étendre]{étendre — to hang out (washing); to spread out}

\label{lesson:FR-C526-etendre}



\begin{quote}
Before the new one: say the French for washing machine, then the French for to spin.
\end{quote}

\subsection*{You'll want to know: étendre}
\textbf{étendre} — \textquotedblleft{}to hang out (washing); to spread out\textquotedblright{}.

\begin{quote}
Elle étend le linge sur le balcon.
\end{quote}

\begin{quote}
\emph{She hangs the washing out on the balcony.}
\end{quote}

\emph{S'étendre} — to lie down, to stretch out.

\begin{cousinweb}
\emph{étendre} is inherited, through Old French \emph{estendre}, from Latin \emph{extendere} \textquotedblleft{}to stretch out\textquotedblright{} (EtymDB), built on \emph{tendere} \textquotedblleft{}to stretch\textquotedblright{} (Lewis \& Short).
\end{cousinweb}

\begin{grammarlens}[title={What you've built}]
One more word for telling, explaining and describing at B1.
\end{grammarlens}

\subsection*{Your turn}
\begin{itemize}
\raggedright
  \item \emph{Say it:} \emph{étendre}
  \item \emph{Say it:} \emph{étendre}, once more
  \item \emph{Say it:} a sentence of your own with \emph{étendre}
  \item \emph{Recall:} say the French for washing machine, then the French for to spin, then say \emph{étendre} again
\end{itemize}

\begin{tcolorbox}[breakable,colback=teal!4,colframe=teal!35!black,title={Before you move on}]
What does étendre mean? (\textquotedblleft{}To hang out (washing); to spread out\textquotedblright{}.) Say it once more.
\end{tcolorbox}
\section[la pince à linge]{la pince à linge — clothes peg}

\label{lesson:FR-C526-pince-a-linge}



\begin{quote}
Before the new one: say the French for to spin, then the French for to hang out.
\end{quote}

\subsection*{You'll want to know: la pince à linge}
\textbf{la pince à linge} — \textquotedblleft{}clothes peg\textquotedblright{}.

\begin{quote}
Il me faut encore des pinces à linge.
\end{quote}

\begin{quote}
\emph{I still need some clothes pegs.}
\end{quote}

\emph{Une pince} alone is pliers or tongs.

\begin{cousinweb}
\emph{la pince à linge} joins \emph{pince}, built on \emph{pincer} \textquotedblleft{}to pinch\textquotedblright{}, an expressive word, and \emph{linge}, from Latin \emph{līneus} \textquotedblleft{}made of linen\textquotedblright{}, from \emph{līnum} \textquotedblleft{}flax\textquotedblright{} (Lewis \& Short).
\end{cousinweb}

\begin{grammarlens}[title={What you've built}]
One more word for telling, explaining and describing at B1.
\end{grammarlens}

\subsection*{Your turn}
\begin{itemize}
\raggedright
  \item \emph{Say it:} \emph{la pince à linge}
  \item \emph{Say it:} \emph{la pince à linge}, once more
  \item \emph{Say it:} a sentence of your own with \emph{la pince à linge}
  \item \emph{Recall:} say the French for to spin, then the French for to hang out, then say \emph{la pince à linge} again
\end{itemize}

\begin{tcolorbox}[breakable,colback=teal!4,colframe=teal!35!black,title={Before you move on}]
What does la pince à linge mean? (\textquotedblleft{}Clothes peg\textquotedblright{}.) Say it once more.
\end{tcolorbox}
\section[le cintre]{le cintre — coat hanger}

\label{lesson:FR-C526-cintre}



\begin{quote}
Before the new one: say the French for to hang out, then the French for clothes peg.
\end{quote}

\subsection*{You'll want to know: le cintre}
\textbf{le cintre} — \textquotedblleft{}coat hanger\textquotedblright{}.

\begin{quote}
Je mets ma chemise sur un cintre.
\end{quote}

\begin{quote}
\emph{I put my shirt on a hanger.}
\end{quote}

In Québec \emph{le support}.

\begin{cousinweb}
\emph{le cintre} is built on \emph{cintrer} \textquotedblleft{}to curve\textquotedblright{} (EtymDB), from Vulgar Latin \emph{cinctūrāre}, from Latin \emph{cinctūra} \textquotedblleft{}a girdle\textquotedblright{}, from \emph{cingere} \textquotedblleft{}to gird\textquotedblright{} (Lewis \& Short).
\end{cousinweb}

\begin{grammarlens}[title={What you've built}]
One more word for telling, explaining and describing at B1.
\end{grammarlens}

\subsection*{Your turn}
\begin{itemize}
\raggedright
  \item \emph{Say it:} \emph{le cintre}
  \item \emph{Say it:} \emph{le cintre}, once more
  \item \emph{Say it:} a sentence of your own with \emph{le cintre}
  \item \emph{Recall:} say the French for to hang out, then the French for clothes peg, then say \emph{le cintre} again
\end{itemize}

\begin{tcolorbox}[breakable,colback=teal!4,colframe=teal!35!black,title={Before you move on}]
What does le cintre mean? (\textquotedblleft{}Coat hanger\textquotedblright{}.) Say it once more.
\end{tcolorbox}
